\section{High-Dimensional Data Visualization}

High-dimensional data has \term{many attributes} (columns / dimensions / features /
variables) per item. A working definition from the lecture: data are
\term{high-dimensional} when our traditional visual encodings for tabular data
\emph{break} (Yang et al., 2003) -- typically already from dozens of dimensions
on. Structured data has named, interpretable attributes (e.g.\ demographic
indicators, gene expression); unstructured data (text, images, audio) is turned
\emph{into} a high-dimensional feature matrix via bag-of-words /
document-term matrices, flattened pixels (MNIST $\to$ $784$-dim vector),
word embeddings (word2vec) or learned latent codes (autoencoders, CNNs).
Abstractly everything is a flat $n \times m$ table $X$ ($n$ items $\times$ $m$
attributes), only $m$ is very large and the attributes may not be interpretable.

\begin{keybox}[Two fundamental strategies]
\begin{itemize}
\item \term{Feature selection}: pick a \emph{subset} of the existing dimensions,
      \emph{without} transforming them. Axes stay interpretable. Dimensions (or
      dimension pairs) are ranked by a \term{quality metric} / ``cognostic''
      (computed diagnostic), then shown with 1D/2D/nD techniques.
\item \term{Feature extraction} (= dimensionality reduction): \emph{transform}
      existing features into a new, lower-dimensional space (project onto new
      axes). The new axes are usually \emph{not} interpretable, and extraction
      \emph{loses information}.
\item \term{Hybrid}: select fewer meaningful dimensions first (needs domain
      knowledge), then apply dimensionality reduction on them.
\end{itemize}
\end{keybox}

% =====================================================================
\subsection{Curse of Dimensionality (Giraud)}

As the dimension $p$ grows, high-dimensional spaces become \term{vast} and data
points get \term{isolated in their immensity}. Three intertwined effects:

\begin{itemize}
\item \term{Distance concentration}: the notion of ``nearest neighbor'' vanishes.
      For $n$ points sampled uniformly in $[0,1]^p$, as $p$ increases the
      \emph{minimal} distance between two points increases and \emph{all}
      pairwise distances become almost equal -- every point is roughly equally
      far from every other. The mean square distance grows linearly with $p$,
      $\mathbb{E}\!\left[\lVert X^{(i)}-X^{(j)}\rVert^2\right]=p/6$, while the
      \emph{relative} spread (scaled deviation) shrinks like $p^{-1/2}$. Any
      method based on \term{local averaging} (e.g.\ $k$-Nearest-Neighbors,
      kernel smoothing) therefore fails.
\item \term{Sparsity}: to keep observations within distance $1$ of any query
      point, the required number of samples $n$ grows \emph{more than
      exponentially} with $p$. The lower bound is
      $n \gtrsim \left(\tfrac{p}{2\pi e}\right)^{p/2}\!\sqrt{p\pi}$; e.g.\ for
      $p=100$ already $n \approx 4\cdot10^{39}$, for $p=200$ more than the
      estimated number of particles in the observable universe. Real data is
      hopelessly sparse in high-$p$ space.
\item \term{Vanishing volume}: the volume $V_p(r)$ of a $p$-ball goes to zero
      more than exponentially fast in $p$; for $p\approx20$ the unit ball is
      already almost $0$. Geometric intuitions from 2D/3D are misleading.
\item \term{Cumulating fluctuations}: accumulation of many small fluctuations
      across directions produces a large global fluctuation; rare events
      accumulate and may no longer be rare; computations/optimizations become
      intensive.
\end{itemize}

\begin{pitfall}[``More features is always better'']
The ``blessing'' of measuring thousands of variables is offset by the curse:
\emph{separating signal from noise is in general almost impossible} in
high-dimensional data. More dimensions $\neq$ more usable information.
\end{pitfall}

% =====================================================================
\subsection{Feature Selection vs.\ Feature Extraction (Ghojogh)}

Both reduce dimensionality: a mapping $\mathbf{x}\in\mathbb{R}^d \mapsto
\mathbf{y}\in\mathbb{R}^p$ with $p\le d$.

\begin{itemize}
\item \term{Feature selection}: $p\le d$ and the selected set
      $\{\mathbf{y}^j\}\subseteq\{\mathbf{x}^j\}$ is an \emph{inclusive subset} of
      the original features. Selection criterion improves/maintains accuracy or
      simplifies the model. There are $2^d$ possible subsets, so heuristics are
      used. Two evaluation philosophies:
      \begin{itemize}
      \item \term{Filter methods}: rank/score features \emph{independently} of any
            learning model, then threshold. Based on \term{relevance} (predictive
            power w.r.t.\ target) and \term{redundancy}. Measures: correlation
            criteria (Pearson), mutual information / information gain,
            $\chi^2$ statistics, Markov blanket, consistency-based filter, FCBF
            (symmetrical uncertainty), Interact, mRMR (max relevance, min
            redundancy). Fast, model-agnostic.
      \item \term{Wrapper methods}: \emph{integrate the model} into the subset
            search -- train and test the classifier on each candidate subset and
            use its performance as fitness. Search is NP-hard, so heuristic:
            sequential forward/backward selection (SFS/SBS, also floating
            SFFS/SFBS) or metaheuristics (genetic algorithms, particle swarm).
            More accurate but expensive; risk of overfitting to the model.
      \item \term{Embedded methods}: selection happens \emph{during} model
            training (e.g.\ LASSO / regularization, tree importances). Between
            filter and wrapper in cost.
      \end{itemize}
\item \term{Feature extraction}: the new set $\{\mathbf{y}^j\}$ is \emph{not} a
      subset but a \emph{different space} (projection to new axes); often
      $p\ll d$. Relies on the \term{manifold hypothesis}: data lie on a
      lower-dimensional sub-manifold. Split into \emph{linear} (PCA) vs
      \emph{non-linear} (kernel PCA, MDS, t-SNE, UMAP, ISOMAP) and
      \emph{unsupervised} (uses only variation, e.g.\ PCA) vs \emph{supervised}
      (uses labels, e.g.\ LDA).
\end{itemize}

% =====================================================================
\subsection{Visual Methods for Multivariate Tabular Data}

These work for a handful to a few dozen dimensions, then break down.

\begin{algobox}[Parallel Coordinates]
\textbf{Core idea:} each dimension is a \emph{vertical axis}, placed side by side;
each item is a \emph{polyline} crossing all axes at its values.\\
\textbf{Key properties:} reading correlation between \emph{adjacent} axes ---
\term{positive correlation} $\Rightarrow$ roughly \emph{parallel} line segments;
\term{negative correlation} $\Rightarrow$ \emph{crossing} (X-shaped) segments.
Good for spotting clusters and trends across many attributes.\\
\textbf{Hyperparameters / choices:} \emph{axis ordering} (only adjacent pairs are
directly comparable, so order matters a lot), axis scaling, opacity.\\
\textbf{Scalability:} suffers heavily from \term{overplotting} with many items;
$\sim$dozens of axes at most.\\
\textbf{When to use:} moderate \#dimensions, want to compare per-item profiles
and adjacent-pair correlations.\\
\textbf{Pitfalls:} non-adjacent correlations are invisible; a poor axis order
hides structure; clutter with large $n$.
\end{algobox}

\begin{algobox}[Star / Radar Plots (Glyphs)]
\textbf{Core idea:} a per-item \term{glyph}: axes radiate from a center, the
value on each axis is a spoke length, endpoints connected into a star polygon.\\
\textbf{Key properties:} compact shape signature per item; good for comparing a
\emph{few} items across a \emph{few} attributes; can be tiled as small multiples.\\
\textbf{Hyperparameters:} axis order (shape depends strongly on it), normalization.\\
\textbf{Scalability:} few items, few dimensions; does not scale to many of either.\\
\textbf{When to use:} profile comparison of individual records/categories.\\
\textbf{Pitfalls:} area is misleading (not a meaningful quantity); shape depends
on dimension order; hard to read with many overlapping glyphs.
\end{algobox}

\begin{algobox}[Scatterplot Matrix (SPLOM)]
\textbf{Core idea:} a grid of all pairwise 2D scatterplots ($m$ dimensions $\to$
$m\times m$ matrix, $\binom{m}{2}$ distinct pairs); diagonal often shows the 1D
distribution.\\
\textbf{Key properties:} every pairwise dependency/correlation is visible; color
by class label to inspect separability.\\
\textbf{Scalability:} number of cells grows as $O(m^2)$ -- breaks down around
$\sim$dozens of dimensions; overplotting per cell with large $n$.\\
\textbf{When to use:} explore all pairwise relations; basis for ranking cells via
\term{scagnostics} / \term{class consistency}.\\
\textbf{Pitfalls:} quadratic blow-up; only \emph{pairwise} (no higher-order)
structure; too small cells become unreadable.
\end{algobox}

\begin{keybox}[Other tabular techniques]
\begin{itemize}
\item \term{Table lens / heatmaps}: the $n\times m$ table itself shown as a color
      matrix (cell value $\to$ color). Scales to many items and dimensions; rows
      and columns can be reordered (seriation) to reveal block structure. Used
      e.g.\ for gene expression and demographic indicators.
\item \term{Glyphs} (general): map data attributes to visual channels of a small
      mark (position, color, shape, size, tilt). Star/radar plots are one family.
\item \term{Star Coordinates} (Kandogan): items as points; each dimension is a
      unit vector from the origin, the point is $\sum_i u_i (d_{ji}-\min_i)$.
      Interactive: scaling an axis $=$ feature emphasis, rotating $=$ change
      correlation, switching off $=$ feature selection. \emph{Strongly depends on
      dimension order.}
\end{itemize}
\end{keybox}

% =====================================================================
\subsection{Ranking Views: Scagnostics \& Class Consistency}

These are \term{quality metrics} for feature \emph{selection}: rank scatterplots
(dimension pairs) so the analyst sees the most ``interesting'' ones first.

\begin{algobox}[Scagnostics (Wilkinson et al., 2005)]
\textbf{Core idea:} \emph{scatterplot diagnostics} -- a set of measures computed
\emph{per scatterplot} to characterize its shape, used to \term{rank} the cells
of a SPLOM.\\
\textbf{The nine measures:}
\begin{itemize}
\item \term{Outlying} -- presence of outliers; via minimum spanning tree (MST):
      $c_{\text{outlying}}=\text{length}(T_{\text{outliers}})/\text{length}(T)$,
      outliers $=$ degree-1 vertices with edge weight $>1.5\,\text{IQR}$.
\item \term{Skewed} -- skewness of the edge-length distribution.
\item \term{Clumpy} -- points form tight clumps/clusters.
\item \term{Sparse} -- points spread thinly / low density.
\item \term{Striated} -- parallel strips/bands of points.
\item \term{Convex} -- how convex the cloud is:
      $c_{\text{convex}}=\text{area}(A)/\text{area}(H)$, $A$ $=$ alpha hull,
      $H$ $=$ convex hull.
\item \term{Skinny} -- elongated, thin shape.
\item \term{Stringy} -- thin, string-like (path-like structure in the MST).
\item \term{Monotonic} -- monotone trend (rank correlation).
\end{itemize}
\textbf{When to use:} too many dimension pairs to inspect by hand; surface plots
with outliers / clusters / trends automatically.\\
\textbf{Pitfalls:} purely shape-based (ignores labels); a measure can be fooled
by sampling/binning; not all ``interesting'' plots score high.
\end{algobox}

\begin{algobox}[Class Consistency (Sips et al., 2009)]
\textbf{Core idea:} for \emph{labeled} data, rank projections by how well classes
are \emph{visually separated}.\\
\textbf{How:} \term{distance-consistency} (DSC) -- the \emph{ratio of data points
whose closest class centroid is their own class's centroid}. High DSC $=$ classes
map to visually separable regions; low DSC $=$ classes overlap. (Slide example:
DSC$=90$ vs DSC$=49$ scatterplots.)\\
\textbf{When to use:} pick the dimension pair / projection that best separates
known classes (e.g.\ for visual inspection of classification).\\
\textbf{Pitfalls:} needs labels; centroid-based, so non-convex / multi-modal
classes can be misjudged.
\end{algobox}

\noindent Related: \term{rank-by-feature} (Seo \& Shneiderman) ranks by 1D
criteria (normality, uniformity/entropy, \#outliers, \#unique values) or 2D
criteria (correlation coefficient, regression error, density);
\term{parallel coordinates in Hough space} (Tatu et al.) -- good dimension pairs
have fewer, well-defined clusters in Hough space.

% =====================================================================
\subsection{Dimensionality Reduction Fact-Sheets}

\begin{algobox}[PCA -- Principal Component Analysis]
\textbf{Core idea:} \term{linear} projection that \term{maximizes variance} along
each new axis, equivalently \term{minimizes the reconstruction (reprojection)
error} (squared distance between original and projected points).\\
\textbf{How:} operate on \emph{normalized} data (centered to mean, scaled to unit
variance). Principal components $=$ \term{eigenvectors of the covariance matrix}
$S$ (solving $S\mathbf{u}=\lambda\mathbf{u}$). The \emph{first} PC is the
eigenvector with the \emph{largest eigenvalue} (direction of max variance); PCs
are \term{mutually orthogonal} and form a new coordinate system.\\
\textbf{Properties:} \term{deterministic}; global, linear; preserves large
pairwise distances; ``score plot'' $=$ projection onto first 2 PCs (proximity
indicates similarity); \term{cannot capture non-linear structure} (fails on the
swiss roll). Often only a small fraction of variance is kept (MNIST PCA $<20\%$).\\
\textbf{Hyperparameters:} number of components $k$ to keep (scree plot,
cumulative explained-variance threshold).\\
\textbf{Complexity:} eigendecomposition of $d\times d$ covariance
($O(d^3)$, or SVD); fast and scalable.\\
\textbf{When to use:} linear correlation structure, need reproducibility,
interpretable variance, preprocessing before t-SNE/UMAP.\\
\textbf{Pitfalls:} only linear; sensitive to scaling (\emph{must} standardize);
2D projection can be a poor approximation if variance is spread over many PCs.
\end{algobox}

\begin{exambox}[Sketching a 1D PCA projection]
A classic exam task: given a 2D point cloud elongated along a diagonal, draw the
\term{first principal component} (the long axis of the cloud, through the mean,
direction of maximum variance) and the second PC orthogonal to it. To get the
\term{1D projection}, drop a \emph{perpendicular} from each point onto PC1 and
read off its position along that line.
\begin{center}
\begin{tikzpicture}[scale=1.0,>={Stealth[round]}]
  % PC1 line direction endpoints used for the perpendicular projection
  \coordinate (O) at (0,0);
  \coordinate (P1) at (1.85,1.46);
  % point cloud along a diagonal; foot = perpendicular projection onto PC1
  \foreach \x/\y in {-1.6/-1.3,-1.1/-0.7,-0.8/-1.0,-0.4/-0.1,-0.2/-0.5,
                     0.1/0.3,0.3/-0.1,0.6/0.7,0.9/0.4,1.3/1.1,
                     1.6/0.9,-0.6/-0.4,1.0/0.9}{
    \fill[tuwblue!55] (\x,\y) circle (1.6pt);
    \draw[dotted,tuwdark!70] (\x,\y) -- ($(O)!(\x,\y)!(P1)$);
    \fill[pitred] ($(O)!(\x,\y)!(P1)$) circle (1.1pt);}
  % PC1 (long axis, slope ~0.79) through the centroid (origin)
  \draw[->,thick,pitred] (0,0) -- (1.85,1.46) node[right]{\footnotesize PC1};
  \draw[pitred] (0,0) -- (-1.85,-1.46);
  % PC2 orthogonal
  \draw[->,thick,keygreen] (0,0) -- (-0.63,0.80) node[above left]{\footnotesize PC2};
  % centroid marker
  \fill[tuwdark] (0,0) circle (1.4pt) node[below right,font=\scriptsize]{mean};
  \node[tuwdark,font=\scriptsize,align=center] at (2.6,-0.55)
        {feet of $\perp$\\$=$ 1D scores};
\end{tikzpicture}
\end{center}
\sketchnote{Ellipse-shaped cloud; arrow PC1 along the long axis through the
centroid; short orthogonal arrow PC2; for each point a thin line dropped at a
right angle onto PC1; the feet of those perpendiculars are the 1D PCA scores.}
\end{exambox}

\begin{algobox}[MDS -- Multidimensional Scaling]
\textbf{Core idea:} find a low-dimensional embedding that \term{best preserves the
pairwise distances} $\delta_{i,j}$ of the high-dimensional data.\\
\textbf{Cost (stress):}
$\displaystyle \min_{\mathbf{x}_1,\dots,\mathbf{x}_n}\sum_{i<j}\bigl(\lVert
\mathbf{x}_i-\mathbf{x}_j\rVert-\delta_{i,j}\bigr)^2$.\\
\textbf{Variants:} \term{classical / metric} MDS preserves actual distance values
(classical MDS on Euclidean distances $\approx$ PCA); \term{non-metric} MDS
preserves only the \emph{rank order} of distances. Distances can be any metric
(e.g.\ cosine similarity for documents). \term{ISOMAP} $=$ MDS on \term{geodesic}
(manifold) distances instead of Euclidean -- a non-linear variant that unrolls
manifolds like the swiss roll.\\
\textbf{Properties:} preserves global distance structure; like PCA tends to keep
\emph{dissimilar} points far apart.\\
\textbf{Hyperparameters:} target dimension, metric, metric vs non-metric.\\
\textbf{Complexity:} works on the $n\times n$ distance matrix $\Rightarrow$
$O(n^2)$ memory/time, heavier for large $n$.\\
\textbf{When to use:} you have a meaningful (dis)similarity matrix and care about
preserving overall distances.\\
\textbf{Pitfalls:} struggles to keep \emph{very similar} points close (crowding);
$O(n^2)$ does not scale to huge $n$.
\end{algobox}

\begin{algobox}[t-SNE -- t-Distributed Stochastic Neighbor Embedding]
\textbf{Core idea:} preserve \term{local neighborhoods}, not distances. Model
high-dimensional similarities as conditional probabilities (Gaussian),
$p_{j|i}\propto\exp(-\lVert x_i-x_j\rVert^2/2\sigma_i^2)$, and low-dimensional
similarities with a heavy-tailed \term{Student-t} distribution $q_{j|i}$;
minimize the \term{Kullback--Leibler divergence}
$C=\sum_i \mathrm{KL}(P_i\Vert Q_i)$ by gradient descent.\\
\textbf{Key properties:} the Student-t tail solves the \term{crowding problem}
(room to spread out medium distances). Keeps similar points close;
\term{stochastic / non-deterministic} (different runs differ).\\
\textbf{Hyperparameters:} \term{perplexity} (effective \#neighbors defining the
Gaussian variance; balances local vs global, $\approx 5$--$50$), learning rate
($\epsilon$), iterations; can use Barnes--Hut approximation.\\
\textbf{Complexity:} heavier than PCA/UMAP; iterative optimization
($O(n^2)$ naive, $O(n\log n)$ with Barnes--Hut).\\
\textbf{When to use:} reveal local cluster structure (e.g.\ MNIST digit clusters).\\
\textbf{Pitfalls (heavily examined):} \emph{cluster sizes are NOT meaningful};
\emph{inter-cluster distances are NOT meaningful}; results depend on perplexity
and random seed; can show clusters in random data at low perplexity; preserves
no global geometry.
\end{algobox}

\begin{algobox}[UMAP -- Uniform Manifold Approximation and Projection]
\textbf{Core idea:} \term{manifold / graph-based}. Build a weighted
$k$-nearest-neighbor graph (a fuzzy topological / simplicial representation of the
manifold), then optimize a low-dimensional layout with the closest equivalent
fuzzy topological structure.\\
\textbf{Key properties:} like t-SNE preserves \term{local} neighborhoods but
tends to preserve \term{more global structure}; \term{faster} and scales better;
still \term{stochastic}.\\
\textbf{Hyperparameters:} \term{n\_neighbors} (size of local neighborhood from
which the manifold is learned -- larger $=$ more global view), \term{min\_dist}
(how tightly points may pack in the embedding -- smaller $=$ tighter clumps).\\
\textbf{Complexity:} approximate kNN $\Rightarrow$ scales to large $n$, faster
than t-SNE.\\
\textbf{When to use:} large datasets needing local \emph{and} some global
structure (single-cell, embeddings).\\
\textbf{Pitfalls:} absolute distances/densities still not fully reliable;
hyperparameter-dependent shapes; non-deterministic.
\end{algobox}

\begin{keybox}[Distortions \& the structure--interpretability trade-off]
Any DR projects feature space $\mathcal{D}$ to visual space $\mathcal{M}$ and
introduces \term{distortions}: \term{missing neighbors} (close in feature space
but far in the plot) and \term{false neighbors} (far in feature space but close
in the plot). Trade-off (Liu et al., 2017): \term{manifold learning} (t-SNE,
UMAP) captures the most intrinsic structure but has \emph{non-interpretable}
axes; \term{axis-aligned projection} (feature selection) has fully interpretable
axes but least structure; \term{linear projection} (PCA) sits in between.
Take-away: \emph{feature extraction leads to information loss}.
\end{keybox}

\subsubsection{Comparison Table: PCA vs MDS vs t-SNE vs UMAP}

\begin{center}
\begin{tabular}{@{}p{1.6cm}p{1.2cm}p{1.4cm}p{3.0cm}p{3.2cm}p{2.0cm}@{}}
\toprule
\textbf{Method} & \textbf{Linear?} & \textbf{Determ.?} & \textbf{Preserves}
& \textbf{Main hyperparameters} & \textbf{Scalability}\\
\midrule
\term{PCA}   & yes & yes & global variance / large distances; linear subspace
             & \#components $k$ & high ($O(d^3)$)\\
\term{MDS}   & no\textsuperscript{*} & yes (classical) & global pairwise distances
             (stress) & target dim; metric; (non)metric & low ($O(n^2)$)\\
\term{t-SNE} & no & \textbf{no} & local neighborhoods (cluster sizes \&
             inter-cluster dist.\ \emph{not} meaningful)
             & \textbf{perplexity}, learning rate $\epsilon$, iterations & medium\\
\term{UMAP}  & no & \textbf{no} & local $+$ more global structure (kNN graph /
             manifold) & \textbf{n\_neighbors}, \textbf{min\_dist} & high (fast)\\
\bottomrule
\end{tabular}
\end{center}
{\footnotesize\textsuperscript{*}Classical/metric MDS on Euclidean distances is
linear and coincides with PCA; non-metric MDS and ISOMAP are non-linear.}

\begin{exambox}[MCQ-style: properties \& hyperparameters]
\begin{itemize}
\item \emph{Which method is deterministic (same result every run)?} $\to$ PCA
      (and classical MDS). t-SNE and UMAP are \term{stochastic}.
\item \emph{Whose cluster sizes and inter-cluster distances are meaningless?}
      $\to$ \term{t-SNE} (and largely t-SNE-style layouts).
\item \emph{``Perplexity'' is a hyperparameter of \dots} $\to$ \term{t-SNE}
      (number of neighbors defining the Gaussian variance, balances local/global).
\item \emph{``n\_neighbors'' and ``min\_dist'' belong to \dots} $\to$ \term{UMAP}.
\item \emph{Which is the only linear method?} $\to$ \term{PCA} (and classical MDS).
\item \emph{PCA finds the directions of \dots} $\to$ \term{maximum variance} $=$
      eigenvectors of the covariance matrix; first PC $=$ largest eigenvalue.
\item \emph{Which best preserves global pairwise distances?} $\to$ MDS / PCA;
      \emph{which best preserves local neighborhoods?} $\to$ t-SNE / UMAP.
\end{itemize}
\end{exambox}

\begin{pitfall}[Over-interpreting a 2D projection]
Do not read meaning into t-SNE/UMAP \emph{distances}, cluster \emph{sizes}, or
empty gaps; do not call a t-SNE result reproducible; do not forget to
\emph{standardize} before PCA; remember that a low total explained variance means
the 2D PCA scatterplot is a weak summary. State that projection axes are
\emph{not interpretable} for feature-extraction methods.
\end{pitfall}
